LLM benchmark scores are reported under many different and often unstated decoding settings: greedy for leaderboards, $T{\approx}0.7$ for chat, and sampled pass@1 for code. If scores, or worse, model \emph{rankings}, depend on the decoder, comparisons across papers are fragile. Prior studies of decoding usually look at one task family at a time and do not measure ranking stability. We run a single controlled grid: 11 open-weight instruction-tuned LLMs (7 families, 0.5B--9B parameters), four domains (factual QA, math reasoning, code synthesis, creative writing), and six decoding strategies (greedy, $T{=}0.7$, $T{=}1.0$, top-$p$ 0.9, top-$k$ 40, min-$p$ 0.1), with three seeds per sampled strategy, for 132k generations in total. In all three verifiable domains, every sampler lowers accuracy on average relative to greedy decoding. Untruncated $T{=}1.0$ costs 4.3 points on \triviaqa, 5.9 on \gsm, and 6.7 on \humaneval. Top-$p$ and min-$p$ recover about half of this loss. Under an LLM-judge proxy, sampling does not improve creative quality (at most $+0.7\%$), and $T{=}1.0$ lowers it by $7.9\%$. Model rankings are robust: Kendall $\tau$ against the greedy ranking is at least 0.85 in every verifiable domain, and $\tau$ falls to 0.71 only for creative writing at $T{=}1.0$. Individual models can still move by 8--15 points. Code is the most seed-sensitive domain: 27\% of problems flip between pass and fail across seeds. We recommend greedy or min-$p$/top-$p$ decoding for evaluation, explicit reporting of decoding settings, and averaging over at least three seeds when sampling.
